\chapter*{Preface}
\addcontentsline{toc}{chapter}{Preface}

Welcome. If you are holding this book, you are probably about to study
algorithms seriously for the first time --- maybe for a course, maybe for an
interview, maybe just because something in you wants to know how computers
actually \emph{think} their way through a problem. Whatever brought you here, I
am glad you came, and I want to make you a promise on the very first page: I am
going to slow down.

Most algorithms books are written by experts, for people who are on their way
to becoming experts. They are precise, they are correct, and they are often
genuinely beautiful. But they move fast. They state a theorem, prove it in five
tight lines, and trust you to fill in the ten lines of intuition that the author
worked out years ago and no longer remembers needing. If you already see it, the
brevity is a gift. If you do not yet see it, the same brevity feels like a wall.

This book is the conversation that happens \emph{before} that wall. It is the
patient friend who sits next to you, draws the picture on a napkin, walks you
through a tiny numeric example by hand, points out the exact place where
everyone gets confused, and only \emph{then} writes down the formal statement ---
at which point the formal statement feels obvious, almost inevitable, because you
already understand the thing it is describing.

\chapterline

\section*{Who this book is for}

You should read this book if any of the following sound like you.

\begin{itemize}
  \item You are about to take an algorithms course, or you are in one right now
        and the textbook feels like it is written in a language you almost speak.
  \item You have done some programming --- you can write a loop, you know what an
        array is, you have used a list or a dictionary --- but nobody ever sat you
        down and explained \emph{why} one way of solving a problem is a thousand
        times faster than another.
  \item You have heard words like ``$O(n \log n)$'', ``recursion'', ``dynamic
        programming'', and ``NP-complete'' and you can repeat them, but you could
        not yet derive them, defend them, or teach them to someone else.
  \item You tried to read CLRS --- \emph{Introduction to Algorithms} by Cormen,
        Leiserson, Rivest, and Stein --- and you bounced off it. That is not a
        verdict on you. CLRS is a magnificent reference and a demanding first
        read. This book exists to make your \emph{second} encounter with it feel
        like a homecoming.
\end{itemize}

You do \emph{not} need to be a strong mathematician. You need to be comfortable
with high-school algebra, willing to count things carefully, and patient enough
to trace an example by hand even when you ``sort of'' see the answer. Every more
advanced tool we need --- summations, logarithms, a little probability, proof by
induction --- is built up gently right where we use it, and collected in the
appendices for reference.

\begin{aestheticbox}
Why does any of this matter once you leave the classroom? Because the difference
between a good algorithm and a careless one is the difference between a search
that returns in 50 milliseconds and one that times out; between a map app that
routes you across a country instantly and one that spins forever; between a
system that scales to a billion users and one that falls over at a thousand.
Algorithms are where good intentions meet the cold arithmetic of \emph{how
many steps}. Learning to count those steps --- and to bring the count down --- is
one of the most durable, transferable skills in all of computing. Hardware gets
faster every year; a better algorithm can make something possible that no amount
of hardware ever would.
\end{aestheticbox}

\section*{The one idea behind the whole book: intuition, then rigor}

Here is the single principle that shapes every chapter. \textbf{First we build a
mental picture you can trust. Then we make it precise.} In that order, never the
reverse.

Why that order? Because precision without a picture is memorization, and
memorization evaporates. If I hand you the recurrence
$T(n) = 2T(n/2) + \Theta(n)$ and tell you it solves to $\Theta(n \log n)$, you
can nod and copy it down, and forget it by Thursday. But if you have \emph{seen}
a list of eight numbers split into two halves, then four, then eight singletons,
and you have counted with your own finger that the splitting takes about $\log_2
8 = 3$ levels and each level does about $n$ units of work --- then the recurrence
is no longer an arbitrary formula to memorize; it is a concise description of a
picture you already understand.

We will draw exactly this kind of picture --- a \emph{recursion tree} --- many
times in the book. Look at one now, even before you can read every
symbol on it. Each node is a piece of the problem; an arrow splits a piece into
the two smaller pieces it asks for help with; the number inside a node is how big
that piece is.

\begin{center}
\begin{tikzpicture}[
    every node/.style={font=\scriptsize},
    pnode/.style={draw, rounded corners, minimum width=0.7cm, minimum height=0.55cm, fill=sourceblue!10},
    lvl/.style={font=\scriptsize\itshape, structurepurple},
    grow=down, level distance=1.05cm,
    level 1/.style={sibling distance=4.6cm},
    level 2/.style={sibling distance=2.3cm},
    level 3/.style={sibling distance=1.15cm}
  ]
  \node[pnode] (root) {$8$}
    child {node[pnode] {$4$}
      child {node[pnode] {$2$}
        child {node[pnode] {$1$}}
        child {node[pnode] {$1$}}}
      child {node[pnode] {$2$}
        child {node[pnode] {$1$}}
        child {node[pnode] {$1$}}}}
    child {node[pnode] {$4$}
      child {node[pnode] {$2$}
        child {node[pnode] {$1$}}
        child {node[pnode] {$1$}}}
      child {node[pnode] {$2$}
        child {node[pnode] {$1$}}
        child {node[pnode] {$1$}}}};
  % level annotations on the right
  \node[lvl, anchor=west] at (4.9, 0)     {level 0: work $\approx n$};
  \node[lvl, anchor=west] at (4.9, -1.05) {level 1: work $\approx n$};
  \node[lvl, anchor=west] at (4.9, -2.10) {level 2: work $\approx n$};
  \node[lvl, anchor=west] at (4.9, -3.15) {level 3: work $\approx n$};
\end{tikzpicture}
\end{center}

The tree has $\log_2 8 = 3$ levels of splitting beneath the root, so four rows in
all. Look across any single row and add up the piece sizes: $8$ on top; $4 + 4 =
8$ next; $2+2+2+2 = 8$; finally eight $1$s. Every row sums to $n = 8$, which is
why each level does about $n$ units of work. Four levels, each costing about $n$,
gives a total of about $n \cdot \log_2 n$ --- and there, drawn rather than
asserted, is the $\Theta(n \log n)$ the recurrence promised.

\begin{intuitionbox}
Think of every algorithm in this book as a \emph{story with a shape}. A loop is
a story where the same small thing happens again and again until a finish line
is crossed. Recursion is a story where a big problem politely asks two smaller
copies of itself for help. A greedy algorithm is a story where you grab the best
thing in front of you and never look back. Before we ever count operations or
prove a bound, we will get the \emph{shape of the story} clear in your head. The
math is what we use to check that the story actually ends well, and to find out
how long it takes.
\end{intuitionbox}

So a typical section moves through four gears:

\begin{enumerate}
  \item \textbf{A picture.} The core mental image, in plain language, often with
        a diagram. ``Imagine the cards in your left hand are already sorted\dots''
  \item \textbf{A tiny worked example.} Real numbers, traced step by step, so you
        can watch the idea move.
  \item \textbf{The pseudocode.} Now that you know what it \emph{does}, here is
        exactly \emph{how} it does it, written carefully.
  \item \textbf{The rigor.} A loop invariant or an induction that proves it is
        correct, and an analysis that derives --- not asserts --- how fast it is.
\end{enumerate}

We will not skip the fourth gear. This is a book that respects proofs; ``it
seems to work'' is never the end of a discussion here. But we earn the proof by
understanding first, so that when the formal statement arrives it reads as a
confirmation of something you already believe, not as a hurdle.

\section*{The colored boxes: a visual language}

Throughout the book you will meet recurring colored boxes. They are not
decoration; each one is a \emph{type of thinking}, and learning to recognize them
on sight will make your reading faster and your studying sharper. Here is the
full vocabulary, with one live example of each so you know exactly what to expect.

\begin{focusbox}
A \textbf{Focus} box holds the one or two things you must carry out of a section.
If you remember nothing else, remember what is in here. Example: \emph{Insertion
sort is fast when the data is nearly sorted and slow when it is reversed --- its
cost is governed by the number of out-of-order pairs, not just the length.}
\end{focusbox}

\begin{intuitionbox}
An \textbf{Intuition} box is the napkin drawing --- the mental model in everyday
words, before any formality. This is where I try to make you \emph{see} the idea.
If a chapter is a visit to new territory, the intuition box is the view from the overlook.
\end{intuitionbox}

\begin{formalbox}
A \textbf{Formal} box states something precisely: a definition, an invariant, a
theorem. After the intuition has done its job, this is the version you could quote
on an exam or cite in a proof. Example: \emph{A sorting algorithm is}
\emph{stable} \emph{if, whenever two elements have equal keys, their relative
order in the output matches their relative order in the input.}
\end{formalbox}

\begin{trapbox}
A \textbf{Common trap} box names a mistake that almost everyone makes, so that you
can make it on purpose, here, safely --- instead of by accident on a test or in
production. Example: \emph{Off-by-one errors cluster around loop bounds. Writing} $i < n$
\emph{when you meant} $i \le n$ \emph{is the single most common bug in all of the
pseudocode you will write this year.}
\end{trapbox}

\begin{walkbox}
A \textbf{Walking question} is a prompt to close the book and think. Genuinely ---
go for a walk, do the dishes, let it simmer. The best learning happens slightly
after you stop reading. Example: \emph{If you could only sort by repeatedly
swapping adjacent elements, what arrangement of} $n$ \emph{numbers would be the
most expensive to sort, and exactly how many swaps would it take?}
\end{walkbox}

\begin{marginnotebox}
A \textbf{History and connection} box tells you where an idea came from --- and
\emph{points directly to the matching place in CLRS} so you always know how this
book maps onto the standard reference. Example: \emph{Quicksort was invented by
Tony Hoare in 1959 while he was a visiting student in Moscow, trying to sort
words for machine translation. We treat it in Chapter~7, exactly mirroring CLRS
Chapter~7.}
\end{marginnotebox}

\begin{aestheticbox}
A \textbf{Why this matters} box connects the idea to the real world --- to the
systems, products, and engineering problems that rely on it in practice.
\end{aestheticbox}

\begin{complexitybox}
A \textbf{Cost summary} box collects the running times and space usage in one tidy
place after we have derived them, so you can find them again at a glance. Example:
\emph{Best case} $\Theta(n)$; \emph{average and worst case} $\Theta(n^2)$;
\emph{extra space} $\Theta(1)$.
\end{complexitybox}

\begin{labbox}
An \textbf{Algorithm lab} is a hands-on ``trace it yourself'' exercise embedded in
the text. You are handed a concrete input and asked to run the algorithm by hand,
writing down the state at each step. There is no substitute for this. Reading an
algorithm and running an algorithm are as different as reading sheet music and
playing it.
\end{labbox}

\begin{formalbox}
There are a few more boxes you will meet occasionally. A \textbf{Reading note}
gives advice on \emph{how} to read a tricky passage. A \textbf{Pause} box is a
gentle checkpoint inviting you to make sure the last idea landed before moving on.
They work just as their names suggest.
\end{formalbox}

\begin{focusbox}
You do not need to memorize this list. After two or three chapters the colors will
do the work for you: green means ``here is the picture,'' grey means ``here is the
precise truth,'' red means ``careful here,'' orange means ``go think,'' and so on.
The boxes are a rhythm. Let them set your pace.
\end{focusbox}

\section*{Where algorithms came from: from mechanical gears to universal recipes}

Before we write down our first loop, we should ask an honest question: why do we
have a discipline called ``algorithms'' at all? Humans have calculated for thousands
of years. Why didn't our ancestors call long division or navigation tables an
``algorithm''?

The answer is a long, dramatic struggle to separate \emph{the recipe} from \emph{the machine
executing it}.

\begin{historybox}
In ninth-century Baghdad, at the House of Wisdom (\emph{Bayt al-Hikma}), Muhammad ibn Musa
al-Khwarizmi wrote a treatise on calculation that gave Europe its system of positional arithmetic
with Hindu numerals; his Latinized name, \emph{Algoritmi}, became our word \emph{algorithm}. For
centuries after, humanity imagined that performing calculations mechanically required building
dedicated, bespoke machines out of brass gears: Blaise Pascal's 1642 Pascaline to add taxes,
Gottfried Wilhelm Leibniz's 1673 Stepped Reckoner to multiply, and Charles Babbage's 1822 Difference
Engine to crank out tables of polynomials. The deep false start of early computing was confusing the
physical mechanism with the idea itself --- believing that each new mathematical problem demanded
an entirely new configuration of metal teeth.

The conceptual rupture arrived with Ada Lovelace in 1843 and Alan Turing in 1936. Translating Luigi
Menabrea's description of Babbage's Analytical Engine, Lovelace saw that the proposed machine did
not merely compute numbers; it could manipulate abstract symbols according to rules, remarking that
``the Analytical Engine weaves algebraical patterns just as the Jacquard-loom weaves flowers and
leaves.'' In Note~G, she published the first true algorithm, a step-by-step trace to calculate
Bernoulli numbers. Nearly a century later, with no electronic computers yet in existence, Alan
Turing imagined a lone clerk manipulating discrete marks on paper tape according to a finite
rulebook. Turing proved that a single Universal Machine could simulate any specialized calculating
machine simply by reading the description of that machine encoded as instructions. The algorithm was
liberated from brass: software was born.

Uncoupling the algorithm from physical hardware unlocked the modern world: a single slab of
silicon can route global communications, render three-dimensional geometry, and sequence genomes
purely by changing the symbolic recipe it runs. Yet this abstraction introduces its own honest
trap. Mathematical algorithm analysis ignores hardware realities: it treats every memory access as
costing the exact same single instant ($\Theta(1)$), ignoring that reading from an on-chip cache
takes a nanosecond while fetching from main memory takes a hundred. An algorithm with a
theoretically superior step count can crawl on real silicon if its memory accesses scatter across
the bus. In this book, we bridge that gap. Chapter~1 examines Euclid's ancient 300~BCE line-measuring
algorithm; Chapter~2 develops loop invariants to prove that programs do what they claim; and
Chapter~3 introduces Donald Knuth's asymptotic language ($\Theta, O, \Omega$) to quantify efficiency
with mathematical honesty.
\end{historybox}

\section*{How the parts map onto CLRS}

This book deliberately follows the structure of CLRS so that the two can be read
side by side. You can use this book as a standalone first course, or you can read
each chapter here and then immediately read the corresponding CLRS chapter to see
the same material at full speed and full rigor. The correspondence is close to
one-to-one. Here is the map.

\begin{longtable}{@{}p{0.07\textwidth}p{0.52\textwidth}p{0.30\textwidth}@{}}
\toprule
\textbf{Ch.} & \textbf{Topic in this book} & \textbf{Matches CLRS} \\
\midrule
\endhead
1  & What is an algorithm?                 & Ch.~1 (the role of algorithms) \\
2  & Getting started: insertion sort       & Ch.~2 \\
3  & Asymptotic notation                   & Ch.~3 \\
4  & Divide and conquer                    & Ch.~4 \\
5  & Probability and randomized algorithms & Ch.~5 \\
6  & Heapsort                              & Ch.~6 \\
7  & Quicksort                             & Ch.~7 \\
8  & Sorting in linear time                & Ch.~8 \\
9  & Medians and order statistics          & Ch.~9 \\
10 & Elementary data structures            & Ch.~10 \\
11 & Hash tables                           & Ch.~11 \\
12 & Binary search trees                   & Ch.~12 \\
13 & Red-black trees                       & Ch.~13 \\
14 & Augmenting data structures            & Ch.~14 \\
15 & Dynamic programming                   & Ch.~15 (Ch.~14 in 4th ed.) \\
16 & Greedy algorithms                     & Ch.~16 (Ch.~15 in 4th ed.) \\
17 & Amortized analysis                    & Ch.~17 (Ch.~16 in 4th ed.) \\
18 & B-trees                               & Ch.~18 \\
19 & Disjoint sets (union-find)            & Ch.~19 / 21 \\
20 & Fibonacci heaps                       & Ch.~19 (3rd ed.) \\
21 & van Emde Boas trees                   & Ch.~20 \\
22 & Elementary graph algorithms           & Ch.~20 / 22 \\
23 & Minimum spanning trees                & Ch.~21 / 23 \\
24 & Single-source shortest paths          & Ch.~22 / 24 \\
25 & All-pairs shortest paths              & Ch.~23 / 25 \\
26 & Maximum flow                          & Ch.~24 / 26 \\
27 & Matrix operations                     & Ch.~28 \\
28 & Linear programming                    & Ch.~29 \\
29 & Polynomials and the FFT               & Ch.~30 \\
30 & Number-theoretic algorithms           & Ch.~31 \\
31 & String matching                       & Ch.~32 \\
32 & Computational geometry                & Ch.~33 \\
33 & NP-completeness                       & Ch.~34 \\
34 & Approximation algorithms              & Ch.~35 \\
35 & Parallel algorithms                   & Ch.~26 (4th ed., multithreading) \\
36 & Machine-learning algorithms           & Ch.~20 (4th ed.) and beyond \\
37 & Modern systems architecture \& memory & Hardware hierarchy, caches, DRAM cycles \\
38 & Compilers: from text to silicon       & Lexing, parsing, SSA, register coloring \\
39 & Quantum computing from the ground up  & Qubits, interference, Grover, Shor \\
\bottomrule
\end{longtable}

A note on the slashes: CLRS has gone through four editions, and the chapter
numbers have drifted between them as topics were added, split, and reordered. Where
a topic sits at different numbers in the third and fourth editions, the table shows
both. Do not let the renumbering rattle you --- the \emph{ideas} are stable, and a
\textbf{History and connection} box at the start of each chapter pins down exactly
where to look in whichever edition you own. Chapters~35 through~39 reach past
the classic CLRS core into material where modern practice and foundational systems live:
parallelism, machine learning, hardware memory hierarchies and DRAM cycles, optimizing
compilers, and quantum computation.

\begin{marginnotebox}
You will notice the early chapters are densely interlinked: insertion sort
(Ch.~2) sets up asymptotic notation (Ch.~3), which is the language used for
everything afterward. If you are reading selectively, treat Chapters~1--4 as the
non-negotiable foundation. After that, the sorting chapters (5--9), the
data-structure chapters (10--14), and the graph chapters (22--26) each form a
fairly self-contained arc you can take in whatever order your course demands.
\end{marginnotebox}

\section*{How to actually use this book}

Reading an algorithms book passively is like watching someone else exercise. You
will enjoy it and learn nothing. Here is how to read this one so that it sticks.

\subsection*{Run the worked examples with a pencil}

Every algorithm comes with at least one fully traced numeric example. \emph{Do
not merely read it.} Cover the worked trace with your hand, predict the next step,
then uncover it and check. When you are wrong --- and you will be, that is the
point --- stop and figure out exactly which assumption in your head was off. That
moment of being wrong and finding out why is where almost all the learning
actually happens.

\subsection*{Do the labs}

The \textbf{Algorithm lab} boxes hand you a fresh input and ask you to be the
computer. Get an index card or a sheet of scratch paper and genuinely write down
the state of the array, the tree, or the queue at each step. It feels slow. It is
slow. It is also the single highest-return activity in the whole book, because it
forces you to hold the precise meaning of every variable in your head at once ---
which is exactly the skill that lets you later \emph{design} algorithms instead of
merely recognizing them.

\begin{labbox}
Here is your very first lab, and it needs no algorithm at all. Take the seven
numbers
\[
  5,\quad 2,\quad 9,\quad 1,\quad 5,\quad 6,\quad 3 .
\]
By hand, sort them into increasing order \emph{using only swaps of adjacent
elements} --- you may only ever exchange two neighbors. Write down each swap you
perform. Count them. Then ask yourself: did you find the smallest possible number
of swaps? How could you be sure? Keep your answer; we will return to exactly this
quantity (it is called the number of \emph{inversions}) when we analyze insertion
sort in Chapter~2.
\end{labbox}

\subsection*{Take the walking questions seriously}

When you hit a \textbf{Walking question}, the best thing you can do is close the
book. The orange boxes are placed at moments where a little independent struggle
will pay off enormously. Your brain keeps working on a good question after you
stop staring at it; some of your best insights will arrive in the shower, on a
walk, or just before sleep. Give them the chance to.

\begin{walkbox}
To break it in right now: you have a phone book with a million names in
alphabetical order, and you want to find ``Ramirez.'' One approach is to read from
the first page until you hit it. Another is to flip to the middle, see whether
``Ramirez'' comes before or after, throw away half the book, and repeat. Roughly
how many flips does the second method need in the worst case? Try to get an actual
number before you read Chapter~1. The gap between ``a million'' and your answer is
the entire reason this field exists.
\end{walkbox}

\subsection*{Read the proofs, but read them gently}

When you reach a loop invariant or an induction, do not brace yourself as if a
storm is coming. A correctness proof in this book is just a careful retelling of
the story you already understand, structured so that no case can sneak through
unhandled. Read it once for the shape of the argument --- ``establish it is true at
the start, show each step preserves it, conclude it holds at the end'' --- and then
once more for the details. If a proof ever feels like it fell from the sky, that is
a bug in my writing, not a deficiency in you; the intuition box right before it is
where I owed you the picture.

\section*{A note on pseudocode}

Throughout the book, algorithms are written in pseudocode --- not in any real
programming language. This is on purpose. Pseudocode lets us express \emph{the
idea} without the clutter of semicolons, type declarations, and library calls that
differ from language to language. Our conventions are deliberately close to those
of CLRS, so that moving between the two books costs you nothing. Here they are, all
in one place.

\begin{itemize}
  \item \textbf{Indentation shows block structure.} The body of a loop or an
        \texttt{if} is the indented text beneath it. There are no curly braces and
        no \texttt{end} keywords cluttering the page.
  \item \textbf{Assignment uses ``$\leftarrow$''} (or sometimes ``='' inside the
        algorithm environments). The statement $x \leftarrow y$ means ``copy the
        value of $y$ into $x$.'' It is not a claim that $x$ and $y$ are equal; it
        is an action. Equality testing, when we need it, is written ``$=$'' inside
        a condition, and we will always make clear from context which we mean.
  \item \textbf{Arrays are indexed from $1$}, the way CLRS does it, unless a chapter
        explicitly says otherwise. So $A[1]$ is the first element and $A[n]$ is the
        last, where $n = A.length$. (If you program in a language that counts from
        $0$, just remember to shift by one when you translate; we will flag the
        places where this bites.)
  \item \textbf{$A[i \mathbin{..} j]$ denotes a slice} --- the elements $A[i],
        A[i+1], \dots, A[j]$. If $i > j$, the slice is empty.
  \item \textbf{Objects have attributes accessed with a dot}, as in $A.length$ or
        $x.key$ or $node.left$. You can read $x.key$ as ``the key field of $x$.''
  \item \textbf{Variables are local} to the procedure they appear in unless we say
        otherwise. Parameters are passed by reference for arrays and objects (so a
        procedure can rearrange an array in place) and by value for plain numbers.
  \item \textbf{We freely use plain English} for steps whose detail would only
        distract --- ``\emph{let} $q$ \emph{be the index of the median element}'' ---
        and then, if that step is the interesting part, we expand it into its own
        procedure.
\end{itemize}

Every algorithm appears inside a numbered, ruled box like the one below, with its
inputs and outputs stated explicitly at the top. Read the \textbf{Input} and
\textbf{Output} lines first --- they are the contract. They tell you what you are
allowed to assume on the way in and what you are promised on the way out, before
you read a single line of the body.

\begin{algorithm}[H]
\caption{A first taste of pseudocode: linear search}
\label{ch00:linear-search}
\KwIn{An array $A[1 \mathbin{..} n]$ and a target value $v$.}
\KwOut{An index $i$ with $A[i] = v$, or the special value \textsc{nil} if no such index exists.}
\For{$i \leftarrow 1$ \KwTo $n$}{
  \If{$A[i] = v$}{
    \Return $i$ \tcp*{found it; stop early}
  }
}
\Return \textsc{nil} \tcp*{checked everything, $v$ is absent}
\end{algorithm}

Take a moment with that little procedure --- it is the simplest interesting
algorithm there is, and it already shows the whole rhythm of the book. The
\textbf{Input} and \textbf{Output} lines state the contract. The loop walks
through the array one slot at a time. The \texttt{if} checks the current slot. The
early \texttt{return} is a shortcut we have earned: the instant we find $v$ there is no
reason to keep looking. And the final \texttt{return} \textsc{nil} handles the case
everyone forgets --- what happens when the thing simply is not there.

We can picture the search as a single pointer sliding rightward along the array,
pausing at each box just long enough to ask ``are you the one?''

\begin{center}
\begin{tikzpicture}[
    box/.style={draw, minimum width=0.95cm, minimum height=0.95cm, font=\small},
    every node/.style={font=\small}
  ]
  \foreach \val/\idx [count=\i from 0] in {4/1, 8/2, 15/3, 16/4, 23/5, 42/6} {
    \ifnum\idx=4
      \node[box, draw=targetgreen, very thick, fill=targetgreen!12] (b\idx) at (\i*1.0, 0) {\val};
    \else
      \node[box] (b\idx) at (\i*1.0, 0) {\val};
    \fi
    \node[gray] at (\i*1.0, -0.75) {\idx};
  }
  % pointer arrow above index 1
  \draw[-{Stealth[length=2.5mm]}, sourceblue, thick] (0, 1.1) -- (0, 0.55);
  \node[sourceblue] at (0, 1.45) {$i$ starts here};
  \node[targetgreen] at (3*1.0, 1.25) {match for $v = 16$};
  \draw[-{Stealth[length=2.5mm]}, targetgreen, thick] (3*1.0, 0.95) -- (3*1.0, 0.55);
\end{tikzpicture}
\end{center}

The pointer $i$ begins at index $1$ and steps right --- past $4$, past $8$, past
$15$ --- until at index $4$ it finds $16$ and stops. It never even looks at $23$
or $42$. If instead we had searched for a value that is not present, the pointer
would have walked all the way off the right end and the final line would have
returned \textsc{nil}.

\begin{trapbox}
The most common bug in a search like this is not in the loop at all; it is
forgetting the \emph{last} line. If you only return inside the \texttt{if}, then a
search for an absent value falls off the bottom of the procedure and returns
garbage (or crashes). Always ask of any algorithm: \emph{what does it do when it
finds nothing?} The empty case, the absent case, the zero-length case --- these boundary
inputs are where subtle bugs hide, so we will check them deliberately in every chapter.
\end{trapbox}

To see why we will spend a whole book worrying about \emph{how} we search and
sort, look at the shape of even this trivial example. In the worst case ---
$v$ is absent, or sits in the last slot --- linear search inspects all $n$
elements. The walking question above hinted that a cleverer search on
\emph{sorted} data can find the same value by inspecting only about $\log_2 n$
elements. For a million elements that is the difference between a million checks
and about twenty. That gap, compounded across millions of queries, is the divide
between an algorithm that glides and one that grinds to a halt.

\begin{complexitybox}
Linear search (Algorithm~\ref{ch00:linear-search}) inspects at most $n$ elements,
so its worst-case running time is $\Theta(n)$, and it uses only $\Theta(1)$ extra
space. We will make the symbols $\Theta$, $\Oh$, and $\Om$ fully precise in
Chapter~3; for now, read $\Theta(n)$ as ``grows in step with $n$'' and $\Theta(1)$
as ``a fixed amount, no matter how big the input.''
\end{complexitybox}

To make that promised gap concrete --- and to give you a second worked example
before you have even reached Chapter~1 --- let us actually run the ``flip to the
middle'' idea on a small sorted array, by hand, the way every section of this book
will ask you to. This is the method the phone-book walking question was circling.
It has a name, \emph{binary search}, and we will meet it properly in Chapter~2;
here we only want to \emph{see} it move.

\begin{labbox}
Take the sorted array
\[
  A = [\,2,\ 5,\ 8,\ 12,\ 16,\ 23,\ 38,\ 56,\ 72,\ 91\,],
\]
ten elements, indexed $1$ through $10$, and search for $v = 23$. Keep two
bookmarks, \emph{lo} and \emph{hi}, marking the part of the array that could still
contain $v$; they start at $1$ and $10$. Repeat: look at the middle element
$A[\,\lfloor(\textit{lo}+\textit{hi})/2\rfloor\,]$, and if it is too small move
\emph{lo} just past it, if too big move \emph{hi} just before it, and if it is
exactly $v$ stop. Before reading the trace below, run those three steps yourself
on paper and count how many middles you inspect.
\end{labbox}

Here is the trace. At each step the live range $[\textit{lo},\textit{hi}]$ is the
only part still in play; everything outside it has already been ruled out by a
single comparison.

\begin{center}
\begin{tabular}{@{}cccccl@{}}
\toprule
\textbf{Step} & \textit{lo} & \textit{hi} & \textbf{mid} & $A[\text{mid}]$ & \textbf{action} \\
\midrule
1 & 1 & 10 & 5 & 16 & $16 < 23$, so search the right half: $\textit{lo} \leftarrow 6$ \\
2 & 6 & 10 & 8 & 56 & $56 > 23$, so search the left half: $\textit{hi} \leftarrow 7$ \\
3 & 6 & 7  & 6 & 23 & $23 = 23$, found at index $6$; stop \\
\bottomrule
\end{tabular}
\end{center}

Three comparisons, and we are done --- against the ten that linear search would
need in the worst case on this same array. The reason is visible in the
\emph{lo}/\emph{hi} columns: each comparison throws away \emph{half} of whatever
remained. Ten candidates become five, then roughly two, then one. The number of
halvings needed to grind a range of $n$ down to a single element is exactly the
quantity $\log_2 n$ that keeps reappearing.

\begin{center}
\begin{tikzpicture}[
    box/.style={draw, minimum width=0.62cm, minimum height=0.62cm, font=\scriptsize},
    dead/.style={box, fill=gray!18, text=gray!55},
    live/.style={box, fill=sourceblue!10},
    midcell/.style={box, draw=quotientorange, very thick, fill=quotientorange!18},
    every node/.style={font=\scriptsize}
  ]
  % row 1: full range, mid = index 5 (value 16)
  \foreach \val/\i [count=\c from 0] in {2/1,5/2,8/3,12/4,16/5,23/6,38/7,56/8,72/9,91/10} {
    \ifnum\i=5
      \node[midcell] (a\i) at (\c*0.66, 0) {\val};
    \else
      \node[live] (a\i) at (\c*0.66, 0) {\val};
    \fi
  }
  \node[sourceblue] at (-1.55, 0) {step 1:};
  % row 2: live = indices 6..10, mid = index 8 (value 56)
  \foreach \val/\i [count=\c from 0] in {2/1,5/2,8/3,12/4,16/5,23/6,38/7,56/8,72/9,91/10} {
    \ifnum\i>5
      \ifnum\i=8
        \node[midcell] (b\i) at (\c*0.66, -0.9) {\val};
      \else
        \node[live] (b\i) at (\c*0.66, -0.9) {\val};
      \fi
    \else
      \node[dead] (b\i) at (\c*0.66, -0.9) {\val};
    \fi
  }
  \node[sourceblue] at (-1.55, -0.9) {step 2:};
  % row 3: live = indices 6..7, mid = index 6 (value 23) -- the hit
  \foreach \val/\i [count=\c from 0] in {2/1,5/2,8/3,12/4,16/5,23/6,38/7,56/8,72/9,91/10} {
    \ifnum\i=6
      \node[box, draw=targetgreen, very thick, fill=targetgreen!18] (c\i) at (\c*0.66, -1.8) {\val};
    \else
      \ifnum\i=7
        \node[live] (c\i) at (\c*0.66, -1.8) {\val};
      \else
        \node[dead] (c\i) at (\c*0.66, -1.8) {\val};
      \fi
    \fi
  }
  \node[sourceblue] at (-1.55, -1.8) {step 3:};
  \node[targetgreen] at (7.1, -1.8) {found $v=23$};
\end{tikzpicture}
\end{center}

The orange cell is the middle we test on each step; the grey cells have been
discarded, never to be looked at again; the blue cells are still in play. Watch
the blue region collapse from ten cells to five to two to the single green hit.
That collapse \emph{is} the algorithm.

\begin{intuitionbox}
This is the whole reason the book exists, in miniature. Linear search and binary
search solve the \emph{identical} problem --- find $v$ in an array --- and a
beginner might call them ``basically the same, just two loops.'' But one inspects
$n$ cells and the other inspects $\log_2 n$, and on a billion-element array that is
the difference between a billion steps and about thirty. Same problem, same
hardware, a thousand-million-fold difference in work --- bought entirely by a
better \emph{shape of story}. Learning to see and create those shapes is the skill
this entire book is here to build.
\end{intuitionbox}

\begin{complexitybox}
On a sorted array of $n$ elements, binary search halves the live range each
comparison, so it finishes after at most $\lfloor\log_2 n\rfloor + 1$ comparisons:
worst-case time $\Theta(\log n)$, extra space $\Theta(1)$. Linear search, by
contrast, costs $\Theta(n)$ time on the same data and does \emph{not} require the
array to be sorted. Sortedness is the price of admission for the logarithmic
speed --- a recurring bargain we will weigh again and again: pay once to organize
the data, then query it cheaply forever after.
\end{complexitybox}

\section*{A few words of encouragement}

Algorithms have a reputation for being hard, and parts of this subject genuinely
are. But a great deal of the difficulty that students experience is not in the
ideas themselves --- it is in the \emph{pace} at which the ideas are usually
delivered. An algorithm that takes an expert thirty seconds to read took that
expert, years ago, an afternoon of confusion to understand for the first time.
Standard textbooks show you the thirty-second version. This book tries to give you
back the afternoon.

So go slowly. Be willing to be confused for a while; confusion is not failure, it
is the feeling of your understanding being under construction. Trace the examples.
Do the labs. Take the walks. Work through the proofs until every step feels
inevitable. If you do those things, then by the time you open CLRS for real, you will
not be meeting these ideas --- you will be greeting old friends, and reading at the
speed the experts do, because by then you will be one.

Let us begin.

\begin{focusbox}
\textbf{The contract of this book.} You bring patience, a pencil, and a
willingness to be temporarily confused. I bring the picture before the proof, a
worked example for everything, and a clear pointer to CLRS at every turn.
Together we get you ready to read the standard reference not with dread, but with
recognition.
\end{focusbox}

\section*{Exercises}

These warm-ups ask for nothing but careful counting and clear thinking --- no
machinery from later chapters. They exist to get you in the habit of \emph{doing}
rather than only reading.

\begin{exercise}
Go back and finish the first \textbf{Algorithm lab}: sort $5, 2, 9, 1, 5, 6, 3$
into increasing order using only swaps of adjacent elements, writing down every
swap. Count your swaps. Now try to find an ordering of these same seven numbers
that requires \emph{more} adjacent swaps to sort, and one that requires
\emph{fewer}. What does the hardest such ordering look like in general?
\end{exercise}

\begin{exercise}
Answer the phone-book \textbf{Walking question}: for a sorted list of one million
names, how many times can you halve the search range before only one name remains?
Give an exact bound, and explain in one sentence why the answer is roughly
$\log_2 n$ rather than $n$.
\end{exercise}

\begin{exercise}
Trace linear search (Algorithm~\ref{ch00:linear-search}) by hand on the array
$A = [4, 8, 15, 16, 23, 42]$ with target $v = 16$. Write down the value of $i$ and
the result of the comparison $A[i] = v$ at every iteration, and state what the
procedure returns.
\end{exercise}

\begin{exercise}
Run the same trace with $v = 17$ on the same array. How many comparisons does the
procedure perform, and what does it return? Which line of the pseudocode is
responsible for producing a sensible answer in this case?
\end{exercise}

\begin{exercise}
Suppose someone hands you a version of linear search that has its final
\texttt{return} \textsc{nil} line deleted. Describe precisely what goes wrong, and
give a concrete input on which the buggy version behaves differently from the
correct one.
\end{exercise}

\begin{exercise}
In your own words, and without looking back, explain the difference between an
\textbf{Intuition} box and a \textbf{Formal} box, and why this book always presents
them in that order. Then invent a one-sentence intuition and a one-sentence formal
statement for the everyday task of ``looking up a word in a dictionary.''
\end{exercise}

\begin{exercise}
Pick any everyday process you do without thinking --- making tea, getting dressed,
finding your keys --- and write it as pseudocode using the conventions from the
``note on pseudocode'' section: an explicit \textbf{Input} line, an explicit
\textbf{Output} line, at least one loop or \texttt{if}, and a comment on at least
one line. Then identify one
\textbf{Common trap}: an input on which your pseudocode would misbehave.
\end{exercise}

\begin{exercise}
Run the binary-search trace from the \textbf{Algorithm lab} again on the same
array $A = [\,2,5,8,12,16,23,38,56,72,91\,]$, but this time search for $v = 12$
and for $v = 40$. For each, write down the $(\textit{lo},\textit{hi},\text{mid})$
triple at every step, the number of comparisons, and what the search concludes.
Which search ends by reporting ``absent,'' and how does the range tell you so?
\end{exercise}

\begin{exercise}
On the same ten-element array, what is the largest number of comparisons binary
search could ever need, over all possible targets? Find a target that achieves it.
Then check your count against the bound $\lfloor\log_2 n\rfloor + 1$ from the
\textbf{Cost summary}.
\end{exercise}

\begin{exercise}[Challenge]
Linear search inspects up to $n$ elements; the phone-book method inspects about
$\log_2 n$. Make a small table of $n = 10,\ 100,\ 1000,\ 10^6,\ 10^9$ and, for
each, write down both $n$ and $\ceil{\log_2 n}$. Looking at your table, write one
sentence explaining why a programmer might reasonably treat $\log_2 n$ as
``practically constant'' for any input that fits on a real computer.
\end{exercise}
